% Vista editorial de corpus/source/masas/02_ley_general.tex, líneas 363--491: cuerpo exacto salvo el retítulo académico del epígrafe sobre el ciclo de 108 estados. Las líneas 1--263 ya residen exactamente en 72_rev2_nucleo_ley_i.tex y las líneas 264--362 en 72_rev2_algoritmo_universal.tex.
\chapter{Extension systématique de la théorie des masses}
\Needspace{7\baselineskip}
\section{Fondement unifié de la masse}
\Needspace{7\baselineskip}
\subsection{L'objet antérieur à la grandeur}
La théorie commence avant la valeur numérique. Une masse physiquement individualisée n'est pas un nombre réel accompagné du nom d'une particule, mais le terme final d'une biographie mathématique. APP fournit l'espace fini d'incidences et les deux opérations corrélatives ; le TRIT conserve l'orientation temporelle de chaque transition ; le TPK transporte phase, feuillet, retenue, frontière, parcours, mémoire et survie. La Mécanique Dimensionnelle convertit ensuite cette biographie en une section de la droite de masse. L'ordre complet est

\[
\mathrm{APP}\longrightarrow\mathrm{TRIT}\longrightarrow\mathrm{TPK}
\longrightarrow\Gamma^{\rm surv}
\longrightarrow\gamma_{\rm obs}
\longrightarrow\mathscr L_\gamma
\longrightarrow\sigma(\gamma)
\longrightarrow\chi_{\rm mass}
\longrightarrow\mathcal T_{90,120}
\longrightarrow\mathcal L_M.
\]
Chaque flèche retient une information utilisée par la suivante. La cellule ne remplace pas le parcours ; le parcours ne remplace pas l'enregistrement ; une signature ne remplace pas son enregistrement généalogique ; un caractère ne remplace pas la section dimensionnelle sur laquelle il agit. Cette séparation explique pourquoi deux particules peuvent partager charge ou spin tout en possédant des masses différentes : leur biographie enrichie n'est pas la même. Elle explique aussi pourquoi particule et antiparticule ont la même masse : l'involution change l'orientation impaire et conserve la mémoire paire qui descend au quotient de masse.

APP comprend neuf fois neuf positions publiées. Chaque position admet quatre orientations initiales, si bien que le domaine orienté contient

\[
81\cdot4=324
\]
parcours candidats. La classification interne produit 56 familles de parcours et 13 multisections. Ces cardinalités sont simultanées : 81 compte les cellules, 324 les orientations, 56 les quotients par histoire et 13 les types opératoires grossiers. Aucun tableau de particules ne remplace ces domaines.

\Needspace{7\baselineskip}
\subsection{Unité discrète du continu qui soutient la masse}
La droite de masse ne s'adosse pas à cinq théories indépendantes du continu. Elle procède de l'unique état APP--TRIT--TPK dont émergent corrélativement les aspects spectral, solénoïdal, de jauge, cohérent et déterminantal, désignés dans la notation de construction par \(sp,sol,gau,coh,det\). Ces symboles sont des discriminants de lecture, et non des sous-domaines préalables. Si \(\mathcal C_{\rm disc}\) désigne la structure commune, la flèche génératrice prend la forme

\[
\mathfrak E_{\rm cont}:\mathcal C_{\rm disc}
\longrightarrow
\bigl(\mathcal C_\infty;
\mathcal O_{sp},\mathcal O_{sol},\mathcal O_{gau},
\mathcal O_{coh},\mathcal O_{det}\bigr),
\]
où les cinq observations appartiennent à un même objet limite, partagent fibres, orientation, retenue, mémoire, frontière et lecteurs, et sont transportées par les mêmes carrés de raffinement. La masse utilise sa lecture déterminantale pour l'échelle d'action, sa lecture spectrale pour le caractère et les tours, sa lecture de jauge pour les secteurs nuls et chargés, et ses lectures cohérente et solénoïdale pour la propagation et la composition. Aucune n'est introduite en supposant d'emblée que l'état appartient à une classe physique classique.

Cette unité est une condition de typage : séparer l'un de ces aspects en conservant les autres produirait un objet distinct, incapable de soutenir simultanément parcours, action, charge, amplitude et masse. La présente théorie ne redémontre pas le continu complet ; elle en incorpore la construction unitaire comme fondement et vérifie qu'aucune des insertions relatives aux masses ne le dissocie.

\Needspace{7\baselineskip}
\subsection{Échelle absolue et quotients}
Le lecteur déterminantal local a pour forme normale de Smith

\[
\operatorname{SNF}(H_4)=\operatorname{diag}(1,2,2,4),
\]
d'où

\[
|\operatorname{coker}H_4|=16,
\qquad
\Sigma_H=\varphi_{\rm HMT}\alpha_{\rm HMT}^{16},
\qquad
\operatorname{ord}_{10}(\Sigma_H)=-34.
\]
La décennie est obtenue avant la carte SI. Les sections pré et de retour de la droite d'action sont

\[
\hbar_{\rm pre}=H_5\,10^{-34}u_S,
\qquad
\hbar_{\rm ret}=\eta_{\rm ret}^{(\deg)}\,10^{-34}u_S,
\]
et leur quotient bidirectionnel est

\[
R_{\rm act}=\frac{H_5}{\eta_{\rm ret}^{(\deg)}}.
\]
Dans la branche d'horloge, nous choisissons explicitement comme section d'action la section de retour, \(\hbar_{\rm int}:=\hbar_{\rm ret}\). Une horloge déclarée \(t_0\) fournit

\[
\omega_0=\frac{2\pi}{108t_0},
\qquad
\varepsilon_{\rm clk}=\hbar_{\rm ret}\omega_0,
\qquad
m_{\rm clk}=\varepsilon_{\rm clk}c_{\rm int}^{-2}.
\]
Adopter \(u_M^{\rm clk}:=m_{\rm clk}\) comme origine de la carte constitue une normalisation explicite. Cela n'identifie pas automatiquement cette section à une base énergétique ou massique préalablement déclarée. Le caractère produit d'abord l'opérateur de base

\[
\widehat M_X^{(0)}
=\mathfrak e_0m_{\rm clk}\,
\chi_{\rm full}(\sigma_X-\sigma_e,\eta_X-\eta_e),
\]
et le lecteur bidirectionnel agit ensuite. Pour chaque présentation \(r\in\{D,\Omega\}\), soit \(\widehat B_X^r=(\widehat B_X^r)^*\) sur la fibre finie ; dans une réalisation infinie, un domaine commun dense est en outre exigé pour le calcul fonctionnel. Alors

\[
\widehat M_X^{\beta,r}
=R_{\rm act}^{\widehat B_X^r/2}
\widehat M_X^{(0)}
R_{\rm act}^{\widehat B_X^r/2}.
\]
Dans une fibre scalaire, \(\widehat B_X^r=\kappa_XI\), de sorte que

\[
m_X^\beta=m_X^{(0)}R_{\rm act}^{\kappa_X}.
\]
Si les coordonnées de tour sont déjà exprimées relativement à l'électron, on omet \(-\eta_e\). Dans un quotient, la section dimensionnelle commune se simplifie, mais non la différence de lecture bidirectionnelle :

\[
\frac{m_Y^\beta}{m_X^\beta}
=\chi_{\rm full}(\sigma_Y-\sigma_X,\eta_Y-\eta_X)
R_{\rm act}^{\kappa_Y-\kappa_X}.
\]
Ainsi, \(10^{-34}\) fixe l'échelle absolue d'action, d'énergie et de masse ; ce n'est pas un correcteur relatif caché. L'information fine d'une espèce ne peut intervenir que par son caractère, son horloge relative ou le quotient orienté des sections pré et de retour.

\Needspace{7\baselineskip}
\subsection{Quantum interne de la droite de masse induit par le cycle de 108 états}
L'autodualité Planck--CT108 fournit, dans la normalisation interne

\[
\hbar_{\rm int}=c_{\rm int}=t_0=1,
\]
la section

\[
m_{108}^{\rm int}
=0.058177641733144319230789692282953757\ldots .
\]
Ce nombre est un quantum normalisé de la droite interne de masse. Il n'est pas la masse d'une particule, ne porte pas à lui seul une carte MeV ou GeV et ne peut être inséré comme ligne supplémentaire d'espèces. Il sert à vérifier la compatibilité entre la carte d'action, l'horloge CT108 et l'autodualité de la droite de masse. La transition vers une masse nominale exige encore le parcours, la signature, l'opérateur et la section dimensionnelle de l'espèce.

\Needspace{7\baselineskip}
\subsection{Existence du résultat HMT–MD pour les extensions survivantes et critère de réduction scalaire selon le rang de la fibre}
Une extension survivante, un lecteur cellulaire, une signature entière, une famille absolument convergente de coefficients de tour et une section dimensionnelle de masse étant fixés, il existe un résultat HMT--MD correctement typé. Si la fibre possède un unique point fixe compatible avec l'involution, le résultat se réduit à un scalaire canonique. Si la fibre est de rang supérieur, le résultat primaire est l'opérateur spectral sur la fibre ; une masse scalaire physique requiert la fonctionnelle de couplage qui individualise la représentation observée.

Ce théorème ne déprécie pas les évaluations scalaires retrouvées. Il distingue trois affirmations : l'opérateur interne est construit ; l'évaluation d'une signature déclarée est exacte ; la détermination prospective de cette signature constitue un théorème supplémentaire lorsqu'elle ne se déduit pas du lecteur lui-même. La comparaison expérimentale n'est effectuée qu'après avoir fixé laquelle de ces trois affirmations est présentée.

\Needspace{7\baselineskip}
